\chapter{Releases e integração contínua}

\eyebrow{como as versões são construídas e distribuídas}

\vspace{4pt}
O projeto é desenvolvido de forma aberta, com verificações de qualidade
automatizadas e builds multiplataforma. Dois workflows do GitHub Actions
conduzem isso: um verifica cada alteração, e outro produz uma release para cada
tag de versão.

\section{Integração contínua}

O workflow de \textbf{CI} roda a cada push e pull request para o branch principal.
Ele verifica as três linguagens do projeto em paralelo:

\begin{center}
\begin{tabularx}{\linewidth}{@{}l X@{}}
\toprule
\textbf{Job} & \textbf{O que verifica} \\
\midrule
Núcleo Go & \tele{go vet} e a suíte de testes completa (\tele{go test ./...}),
      incluindo os testes de integridade de pareamento e de $\kappa$ de Fleiss. \\
Build Wails & Um \tele{wails build} completo: gera os bindings Go$\leftrightarrow$JS,
      faz a verificação de tipos do front end (\tele{tsc}), o compila (\tele{vite})
      e compila o núcleo Go, validando as três camadas juntas, do mesmo modo que
      uma release. \\
Sidecars Python & Compilação de bytes (sintaxe) mais um lint com \tele{ruff} dos
      sidecars de inferência, conformação, ao vivo e Kinect. \\
\bottomrule
\end{tabularx}
\end{center}

Um selo verde de CI no repositório significa que os três passaram no branch
principal atual.

\section{Releases}

O workflow de \textbf{Release} é disparado quando uma tag de versão que
corresponde a \tele{v*} é enviada. Ele constrói a aplicação nativa em três
runners em paralelo (macOS, Windows e Linux), empacota os pesos do modelo em
cada build e anexa os arquivos a uma GitHub Release com notas geradas
automaticamente.

\begin{center}
\begin{tabularx}{\linewidth}{@{}l l@{}}
\toprule
\textbf{Plataforma} & \textbf{Arquivo da release} \\
\midrule
macOS (universal) & \tele{carcass-macos.tar.gz} \\
Windows (amd64)   & \tele{carcass-windows.zip} \\
Linux (amd64)     & \tele{carcass-linux.tar.gz} \\
\bottomrule
\end{tabularx}
\end{center}

Cada arquivo é autocontido: os pesos treinados (\tele{fat\_binary.pth},
\tele{direct\_class.pth}, \tele{eg\_regression.pth}) e a tabela de consulta
ColorNaming viajam junto com o binário, colocados onde o app os espera em tempo
de execução.

\section{Publicação de uma nova versão}

A publicação de uma release consiste em um único push com tag:

\begin{lstlisting}
git tag -a v0.2.0 -m "Apex v0.2.0 - <summary>"
git push origin v0.2.0
\end{lstlisting}

O workflow executa o restante: os três builds de plataforma, o empacotamento dos
pesos e a release publicada com binários prontos para download. A versão
em \tele{wails.json} deve ser mantida em sincronia com a tag.

\begin{note}[title={Reprodutibilidade do build}]
Como a release usa o mesmo \tele{wails build} executado pela CI a cada commit, um
build com tag é consistente com a verificação de CI: se a CI está verde, a release
compila. Os binários resultantes são executáveis sem ferramentas de build. Os
recursos de modelo requerem um ambiente Python com as bibliotecas científicas.
\end{note}
